\documentclass[11pt]{article}

\usepackage[margin=1in]{geometry}
\usepackage{amsmath, amssymb, amsthm}
\usepackage{microtype}
\usepackage{xcolor}
\usepackage{hyperref}
\usepackage{enumitem}
\usepackage{parskip}

\hypersetup{
  colorlinks=true,
  linkcolor=blue!60!black,
  urlcolor=blue!60!black,
  citecolor=blue!60!black
}

\newcommand{\M}{\mathbb{M}}
\newcommand{\C}{\mathbb{C}}
\newcommand{\Z}{\mathbb{Z}}

\title{\textbf{The Monster Group} $\M$}
\author{}
\date{}

\begin{document}
\maketitle

\vspace{-2em}

The largest of the 26 \emph{sporadic simple groups}. To say what that
means we have to back up.

\section*{What is a sporadic group?}

\subsection*{Groups, briefly}

A \emph{group} is a set with a way to combine two elements (think
``multiplication'') that is associative, has an identity, and has
inverses. Examples:
\begin{itemize}[leftmargin=*, itemsep=2pt]
  \item the integers $\Z$ under addition,
  \item the symmetries of a square (rotations and reflections, 8 of them),
  \item the permutations of $n$ objects --- the symmetric group $S_n$,
        size $n!$.
\end{itemize}
A group is \emph{finite} if it has finitely many elements. The
``shape'' of a finite group is captured by its multiplication table.

\subsection*{Normal subgroups and ``simple''}

Given a group $G$, certain subgroups $N \subseteq G$ are
\emph{normal}: roughly, they sit symmetrically inside $G$ so that you
can ``divide them out'' and form a smaller \emph{quotient group}
$G/N$. Normal subgroups are the analogue, for groups, of factors
of an integer.

A group is \textbf{simple} if it has no nontrivial normal subgroup
(no proper $N$ with $1 < |N| < |G|$). Simple groups are the
``primes'' of group theory: every finite group can be broken apart
into simple pieces, and those pieces are unique (Jordan--H\"older).
So if you want to understand all finite groups, the first job is to
list all the finite simple groups.

\subsection*{The classification}

This list was finished around 2004, after roughly fifty years of work
by hundreds of mathematicians, in tens of thousands of journal pages.
The result --- the \emph{Classification of Finite Simple Groups} ---
says that every finite simple group is one of:

\begin{enumerate}[leftmargin=*, itemsep=2pt]
  \item $\Z/p\Z$ for a prime $p$ (cyclic groups of prime order, the
        only abelian simple groups);
  \item $A_n$ for $n \ge 5$ (the alternating groups: even
        permutations of $n$ objects);
  \item a group of \emph{Lie type} (16 infinite families, built from
        matrix groups over finite fields --- things like
        $\operatorname{PSL}_n(\mathbb{F}_q)$);
  \item or one of \textbf{26 sporadic groups}, which fit into none of
        the above families.
\end{enumerate}

The first three categories are infinite families parameterised by
$p$, $n$, or $(n,q)$ --- regular, predictable, generated by a recipe.
The 26 sporadics are the leftovers: isolated exceptions that the
recipes do not produce. They were each discovered one at a time,
often by a person who suspected a particular group ``ought to exist''
based on combinatorial or geometric constraints and then constructed
it by hand.

\subsection*{Brief timeline (see detailed version below)}

\begin{itemize}[leftmargin=*, itemsep=2pt]
  \item \textbf{1861--1873.} Mathieu finds the first five.
  \item \textbf{1965.} Janko discovers $J_1$ after a 92-year gap.
  \item \textbf{1965--1976.} Twenty more sporadics arrive, mostly
        predicted from involution centralisers.
  \item \textbf{1980.} The Monster and $J_4$ are constructed; the
        list closes at 26.
\end{itemize}

\subsection*{The Happy Family and the Pariahs}

Of the 26 sporadics, 20 of them (including the Monster itself) sit
inside $\M$ as subquotients --- you can find them by taking a
subgroup of $\M$ and then quotienting. Robert Griess named these
twenty the \emph{Happy Family}. The remaining six --- $J_1$, $J_3$,
$J_4$, $\mathit{Ru}$, $\mathit{Ly}$, $\mathit{ON}$ --- have no such
relationship to the Monster and are called the \emph{pariahs}.

That the Monster contains most of the other sporadic groups
\emph{inside} it is part of why it is the natural climax of the
story.

\section*{Detailed timeline: the 26 sporadics in order of discovery}

What follows is the canonical chronology. Where a group was
\emph{predicted} from theoretical constraints before being
\emph{constructed} as a concrete object, both dates are given. The
predictor and constructor are often different people; this was normal
in the 1960s and 1970s.

\subsection*{Era I --- The Mathieu groups (1861--1873)}

\begin{itemize}[leftmargin=*, itemsep=2pt]
  \item \textbf{1861.} \'Emile Mathieu introduces $M_{12}$ and
        $M_{24}$, with $M_{11}$ as a subgroup of $M_{12}$.
  \item \textbf{1873.} Mathieu completes the picture with $M_{22}$
        and $M_{23}$, sitting inside $M_{24}$.
\end{itemize}

These five groups (the \emph{Mathieu groups}) sit on top of the
\emph{Steiner system} $S(5,8,24)$ and the binary \emph{Golay code}.
After Mathieu, the subject went silent for 92 years --- people
assumed there were no more sporadic groups, or did not know how to
look.

\subsection*{Era II --- The modern era opens (1965--1968)}

\begin{itemize}[leftmargin=*, itemsep=2pt]
  \item \textbf{1965.} Zvonimir Janko discovers $J_1$ by analysing a
        hypothetical involution centraliser of the form
        $\Z_2 \times A_5$. $J_1$ is a $7\times 7$ matrix group over
        $\mathbb{F}_{11}$. It is the first new sporadic group in
        almost a century, and it convinces everyone that more must
        exist.
  \item \textbf{1967.} Donald Higman and Charles Sims, while
        attending a conference, construct the
        \emph{Higman--Sims group} $\mathit{HS}$ over a single
        weekend as a permutation group on $100$ points.
  \item \textbf{1968.} Conway computes
        $\operatorname{Aut}(\Lambda_{24})$ and identifies
        $\mathit{Co}_1$, $\mathit{Co}_2$, $\mathit{Co}_3$. Three
        sporadics in one calculation.
  \item \textbf{1968.} Janko predicts a second group $J_2$
        (= $\mathit{HJ}$, the \emph{Hall--Janko group}); Marshall
        Hall constructs it as a rank-$3$ permutation group.
  \item \textbf{1968.} Michio Suzuki introduces the \emph{Suzuki
        sporadic group} $\mathit{Suz}$, inside $\mathit{Co}_1$.
  \item \textbf{1968.} Jack McLaughlin introduces the
        \emph{McLaughlin group} $\mathit{McL}$, related to the Leech
        lattice.
  \item \textbf{1968.} Dieter Held predicts the \emph{Held group}
        $\mathit{He}$, with involution centraliser modelled on
        $M_{24}$. Constructed shortly after by Graham Higman and
        John McKay.
\end{itemize}

\subsection*{Era III --- Janko's third, Fischer's family (1969--1972)}

\begin{itemize}[leftmargin=*, itemsep=2pt]
  \item \textbf{1969.} Janko predicts $J_3$. G.\ Higman and J.\ McKay
        construct it (a third Janko group, unrelated to $J_2$ in any
        nice way).
  \item \textbf{1971.} Bernd Fischer, studying groups generated by
        \emph{$3$-transpositions} (involutions whose pairwise
        products have order $1$, $2$, or $3$), introduces three new
        sporadics: $\mathit{Fi}_{22}$, $\mathit{Fi}_{23}$,
        $\mathit{Fi}_{24}'$. The Fischer groups form a tower:
        $\mathit{Fi}_{22} \subset \mathit{Fi}_{23} \subset
        \mathit{Fi}_{24}'$.
  \item \textbf{1972.} Richard Lyons predicts the \emph{Lyons group}
        $\mathit{Ly}$ from an involution centraliser modelled on
        $\mathit{McL}$. Charles Sims constructs it in $1973$ using a
        $111$-dimensional representation over $\mathbb{F}_5$ ---
        an early major use of computers in sporadic-group
        construction.
\end{itemize}

\subsection*{Era IV --- Monster prediction and the final stretch (1973--1976)}

\begin{itemize}[leftmargin=*, itemsep=2pt]
  \item \textbf{1973.} Arunas Rudvalis predicts the
        \emph{Rudvalis group} $\mathit{Ru}$. Conway and David Wales
        construct it later that year.
  \item \textbf{1973.} \emph{Fischer and Griess independently
        predict the Baby Monster $B$ and the Monster $\M$}, both
        much larger than anything previously contemplated.
  \item \textbf{1975.} Janko predicts $J_4$ from an exotic
        involution centraliser involving $M_{22}$.
  \item \textbf{1975--76.} Michael O'Nan introduces the
        \emph{O'Nan group} $\mathit{ON}$, predicted from a
        hypothetical involution centraliser.
  \item \textbf{1976.} Koichiro Harada and Simon Norton introduce the
        \emph{Harada--Norton group} $\mathit{HN}$, predicted as an
        involution-centraliser quotient inside the Monster.
  \item \textbf{1976.} John Thompson predicts the \emph{Thompson
        group} $\mathit{Th}$ from a hypothetical centraliser
        connected to the exceptional Lie group $E_8(\mathbb{F}_3)$.
        Smith and Thompson construct it shortly afterwards.
\end{itemize}

\subsection*{Era V --- The list closes (1977--1980)}

\begin{itemize}[leftmargin=*, itemsep=2pt]
  \item \textbf{1977.} Jeff Leon and Charles Sims construct the
        Baby Monster $B$ by computer, using a permutation
        representation on $13{,}571{,}955{,}000$ points.
  \item \textbf{1980.} Robert Griess constructs the Monster $\M$ by
        hand, as the automorphism group of his $196{,}884$-dimensional
        Griess algebra. He writes the construction up in
        \emph{The Friendly Giant} ($1982$).
  \item \textbf{1980.} Simon Norton, Richard Parker, John Conway and
        David Benson construct $J_4$ using a $112$-dimensional
        representation over $\mathbb{F}_2$, with substantial
        computer help.
\end{itemize}

That is all 26. After August $1980$, every involution centraliser
that arose in the ongoing classification effort matched a group
already on the list. Nobody has ever produced a credible candidate
for a 27th.

\section*{Size of the Monster}

\[
|\M| = 808\,017\,424\,794\,512\,875\,886\,459\,904\,961\,710\,757\,005\,754\,368\,000\,000\,000
\]

\noindent which is about $8 \times 10^{53}$. The prime factorisation is
\[
|\M| = 2^{46} \cdot 3^{20} \cdot 5^9 \cdot 7^6 \cdot 11^2 \cdot 13^3 \cdot 17 \cdot 19 \cdot 23 \cdot 29 \cdot 31 \cdot 41 \cdot 47 \cdot 59 \cdot 71 .
\]

\noindent The 15 primes appearing here are the \emph{supersingular primes}:
exactly the primes $p$ for which the modular curve $X_0(p)^+$ has genus
zero. This coincidence has no business being true and was an early hint
that something deep was happening.

\section*{History}

\begin{itemize}[leftmargin=*, itemsep=2pt]
  \item \textbf{1973.} Bernd Fischer and Robert Griess independently
        predict its existence from group-theoretic constraints.
  \item \textbf{1980.} Griess constructs $\M$ explicitly as the
        automorphism group of a 196{,}884-dimensional commutative
        non-associative algebra --- the \emph{Griess algebra}.
        His paper is titled ``The Friendly Giant.''
\end{itemize}

\section*{Smallest representations}

\begin{itemize}[leftmargin=*, itemsep=2pt]
  \item Trivial: dimension $1$.
  \item Smallest faithful complex representation: dimension
        $\mathbf{196{,}883}$.
  \item The Griess algebra: dimension $196{,}884 = 196{,}883 + 1$.
\end{itemize}

\section*{Why exactly 26? Architecture of the classification proof}

There is no single short proof. The classification of finite simple
groups is a body of roughly $10{,}000$--$15{,}000$ journal pages
across about $500$ papers by more than $100$ mathematicians, spanning
roughly $1955$ to $2004$. What follows is the \emph{structure}: the
load-bearing theorems and the strategy that ties them together.

\subsection*{Step 1.\quad Reduce to groups with an involution}

\textbf{Feit--Thompson Odd Order Theorem (1963).} Every finite group
of odd order is solvable.

\noindent A solvable group is built from abelian pieces, and the only
abelian simple groups are $\Z/p\Z$. So: \emph{every nonabelian finite
simple group has even order}, and therefore contains an element of
order $2$ (an \emph{involution}). This is the foundation; the proof
alone runs $\approx 255$ pages and was a landmark.

After Feit--Thompson, the search reduces to nonabelian finite simple
groups, all of which have involutions.

\subsection*{Step 2.\quad Use the involution centraliser as a fingerprint}

For $t \in G$ an involution, the \emph{centraliser}
$C_G(t) = \{g \in G : gt = tg\}$ records how $t$ sits inside $G$.

\textbf{Brauer--Fowler Theorem (1955).} For any fixed finite group
$H$, only finitely many finite simple groups $G$ have an involution
$t$ with $C_G(t) \cong H$.

\noindent This is the theorem that makes the entire program possible.
It says the involution centraliser is a near-fingerprint: pin it
down, and the simple group is determined up to finitely many
possibilities. The strategy of the classification is therefore:

\begin{enumerate}[leftmargin=*, itemsep=2pt]
  \item enumerate the possible involution centralisers $H$ that could
        occur in a nonabelian finite simple group;
  \item for each $H$, identify the (finitely many) simple groups it
        produces.
\end{enumerate}

\subsection*{Step 3.\quad The big dichotomy}

A nonabelian finite simple group $G$ falls into one of two regimes,
depending on the structure of $C_G(t)$:

\begin{itemize}[leftmargin=*, itemsep=2pt]
  \item \textbf{Characteristic-$2$ type.} The centraliser of every
        involution has a generalised Fitting subgroup that is a
        $2$-group. These groups behave like Chevalley groups in
        characteristic $2$.
  \item \textbf{Component type.} Some involution centraliser has a
        quasi-simple ``component.'' These groups behave like Chevalley
        groups in odd characteristic, or like alternating groups, or
        like sporadics.
\end{itemize}

The classification proof is, at the top level, a vast case analysis
across this dichotomy. The two halves use different tools:
Aschbacher's \emph{component analysis} and the \emph{signaliser
functor method} (Gorenstein--Walter, Glauberman, Goldschmidt) on the
component side; the \emph{Bender method}, $2$-local analysis, and
weak-closure arguments on the characteristic-$2$ side.

\subsection*{Step 4.\quad For each branch, match against a known list}

In each leaf of the case analysis, one of these things must happen:

\begin{itemize}[leftmargin=*, itemsep=2pt]
  \item $C_G(t)$ matches the involution centraliser of a known
        Chevalley group $\Rightarrow$ $G$ \emph{is} that Chevalley
        group (this is the bulk of the work and rules out enormous
        territory in one stroke);
  \item $C_G(t)$ matches the involution centraliser of an alternating
        group $\Rightarrow$ $G \cong A_n$;
  \item $C_G(t)$ matches the involution centraliser of a known
        sporadic $\Rightarrow$ $G$ is that sporadic;
  \item $C_G(t)$ matches no known group $\Rightarrow$ a new sporadic
        is constructed and added to the list.
\end{itemize}

The last bullet is the only way new sporadics can enter. Each of the
26 was found exactly this way: someone was working through the case
analysis, hit an involution centraliser that did not match any known
group, computed what such a group would have to look like (order,
character table, key subgroups), and then constructed it. After
$J_4$ was constructed in $1980$, every branch terminated in an
already-known group. No further branches remain.

\subsection*{Step 5.\quad Verify the case analysis is exhaustive}

This is the part that is genuinely $\sim 10{,}000$ pages. The
exhaustiveness is established in named subtheorems --- e.g.\
Aschbacher's \emph{classical involutions theorem} (1977), the
\emph{trichotomy theorem}, the \emph{quasithin theorem}.

\subsection*{The gap, and the second-generation proof}

In $1981$ Gorenstein announced the classification was complete. In
$1989$ a gap was found in the \emph{quasithin} case (groups of
characteristic-$2$ type with small $2$-local rank). The gap was not
closed until $2004$, when Aschbacher and Smith published a
$1{,}221$-page two-volume monograph filling it. With that, the
classification is universally accepted.

A separate effort --- the \emph{second-generation proof}, started by
Gorenstein, Lyons and Solomon in the $1990$s --- aims to give a more
unified and streamlined account. It is roughly $13$ volumes long and
still in progress.

\subsection*{Why no 27th?}

For a 27th sporadic group $G$ to exist, every branch of the case
analysis would already have to be wrong: some involution
centraliser $H$ would have to admit a simple group not on the list,
and the branch of the proof that should have ruled this out would
have to be incorrect. The two largest such risks --- the quasithin
case and the uniqueness theorems for the largest sporadics --- have
both been re-examined in detail. The proof is not aesthetically
clean, but it is exhaustive.

This is why ``find a new one'' is not a research task one could pursue
in good faith. It is the same shape of claim as ``find a new prime
between $7$ and $11$.''

\section*{$J_4$: a pariah}

Zvonimir Janko's fourth and final sporadic group, and one of the six
\emph{pariahs} --- the sporadics that do \emph{not} embed inside the
Monster.

\textbf{Order.}
\[
|J_4| \;=\; 2^{21} \cdot 3^3 \cdot 5 \cdot 7 \cdot 11^3 \cdot 23 \cdot 29 \cdot 31 \cdot 37 \cdot 43
\;\approx\; 8.68 \times 10^{19}.
\]

\textbf{Smallest faithful representation.} Dimension $1{,}333$ over
$\mathbb{F}_{11}$ (and dimension $1{,}333$ over $\C$ as well). The
prime $11$ plays a central structural role: $|J_4|$ is divisible by
$11^3$, and the group can only be seen comfortably in characteristic
$11$.

\textbf{History.} Predicted by Janko in 1975 on the basis of a
hypothetical involution centraliser. The actual construction was a
heroic effort by Simon Norton, Richard Parker, John Conway, David
Benson and others, using a $112$-dimensional representation over
$\mathbb{F}_2$ and substantial computer work. It is the last sporadic
group ever discovered: by the time $J_4$ was constructed in 1980,
Griess had constructed the Monster, and the classification was
complete.

\textbf{Why pariah?} The Monster's involution centralisers are
double covers of $B$ and of $2^{1+24}.\mathit{Co}_1$. Neither of
these contains $J_4$, so $J_4$ has no home inside $\M$ --- it is
sporadic in the literal sense of standing apart, even from the rest
of the family.

\section*{The Baby Monster $B$}

The second-largest sporadic group, and the most natural companion to
the Monster.

\textbf{Order.}
\[
|B| \;=\; 2^{41} \cdot 3^{13} \cdot 5^6 \cdot 7^2 \cdot 11 \cdot 13 \cdot 17 \cdot 19 \cdot 23 \cdot 31 \cdot 47
\;\approx\; 4.15 \times 10^{33}.
\]
Roughly $10^{20}$ times smaller than $\M$, but still vastly bigger
than any of the other 24 sporadics.

\textbf{Smallest faithful representation.} Dimension $4371$ over $\C$.
(Compare $196{,}883$ for the Monster.) Like the Monster, $B$ has no
tiny faithful representation, just a much less extreme version of the
same phenomenon.

\textbf{History.} Predicted by Bernd Fischer in the early 1970s ---
around the same time as the Monster --- and constructed shortly after.
Fischer called it the ``baby monster'' as shorthand while he was
hunting for the larger group.

\textbf{Relation to $\M$.} Pick an involution $z \in \M$ of conjugacy
class 2A (one of the two classes of order-2 elements in $\M$). Its
centraliser is a double cover $2.B$, and quotienting by $\langle
z\rangle$ recovers $B$ itself. So the Baby Monster is, almost
literally, ``what you see when you look at $\M$ through a single
involution.''

\textbf{Moonshine for $B$.} Just as the Monster has Monstrous
Moonshine attached to the $j$-function, $B$ has its own
\emph{generalised moonshine} story: its character tables show up in
the $q$-expansions of modular functions associated to particular
elements of $\M$. The Baby Monster's moonshine is a corollary of
the Monster's, but historically it was one of the early clues that
the moonshine phenomenon was structural, not a coincidence
restricted to a single group.

\section*{Monstrous Moonshine}

In 1978, John McKay noticed:
\[
196{,}884 \;=\; 196{,}883 + 1.
\]

\noindent The left-hand side is a coefficient of the modular $j$-function,
\[
j(\tau) \;=\; q^{-1} + 744 + 196{,}884\,q + 21{,}493{,}760\,q^2 + 864{,}299{,}970\,q^3 + \cdots
\]
where $q = e^{2\pi i \tau}$. The right-hand side is a sum of dimensions
of irreducible $\M$-representations. The pattern continues:
\begin{align*}
21{,}493{,}760  &= 1 + 196{,}883 + 21{,}296{,}876 ,\\
864{,}299{,}970 &= 2\cdot 1 + 2\cdot 196{,}883 + 21{,}296{,}876 + 842{,}609{,}326 .
\end{align*}

\noindent Every coefficient of $j$ is a small positive-integer combination
of dimensions of Monster irreducible representations.

Conway and Norton formalised this as the \emph{Monstrous Moonshine}
conjectures: for each $g \in \M$, there is a modular function $T_g$
whose $q$-expansion coefficients are characters of $g$ acting on a
graded $\M$-module $V^{\natural}$.

\textbf{Borcherds} proved this in 1992 (Fields Medal, 1998) by
constructing $V^{\natural}$ --- the \emph{moonshine module} --- as a
vertex operator algebra, using generalised Kac--Moody algebras and
ideas from string theory. The Monster turns out to be the symmetry
group of a particular bosonic-string compactification on the Leech
lattice modulo $\Z_2$.

\section*{Timeline of the classification proof}

The proof of ``there are exactly $26$ sporadic groups'' is really a
proof of the larger theorem that classifies \emph{all} finite simple
groups. Sporadics are the leftovers, so finishing the classification
finishes the count. Here are the milestones that matter, with each
plainly explained.

\begin{itemize}[leftmargin=*, itemsep=4pt]
  \item \textbf{1955 --- Brauer--Fowler.} \emph{Plain version:} if
        you tell me how an order-$2$ element sits inside a finite
        simple group (technically: its centraliser), I can narrow
        the group down to finitely many possibilities. This is what
        turns the classification from ``infinite haystack'' into
        ``finite checklist.''
  \item \textbf{1963 --- Feit--Thompson.} \emph{Plain version:}
        every nonabelian finite simple group has an order-$2$
        element to work with. Without this you cannot even start
        Brauer--Fowler's strategy. The proof is $255$ pages and
        was a milestone in its own right.
  \item \textbf{Late 1960s--early 1970s --- the toolkit.}
        Mathematicians (Bender, Glauberman, Gorenstein, Walter,
        Aschbacher) build the machinery for analysing how
        order-$2$ elements interact: the \emph{signaliser functor
        method}, $2$-\emph{local analysis}, \emph{component
        analysis}. \emph{Plain version:} each is a tool for turning
        a hypothetical centraliser into facts about the rest of the
        group.
  \item \textbf{1972 --- Gorenstein's program.} Daniel Gorenstein
        lectures at the University of Chicago, laying out a $16$-step
        roadmap for completing the classification. Up to this
        point, the field had been a chaotic land grab; afterwards,
        people knew which case to take.
  \item \textbf{1976 --- Aschbacher's classical involutions
        theorem.} \emph{Plain version:} a huge chunk of the
        checklist is closed in one paper. Groups with an involution
        whose centraliser ``looks classical'' are shown to be the
        already-known Lie-type groups. No surprises in that branch.
  \item \textbf{1976--1977 --- the trichotomy theorem.} The
        remaining cases are split into three big buckets, each
        attacked separately. This is the structural skeleton of the
        rest of the proof.
  \item \textbf{1980 --- the last sporadics are built.} Griess
        constructs the Monster. Norton, Parker, Conway and Benson
        construct $J_4$. After this, every open branch of the
        classification case analysis has been resolved by either a
        Lie-type group, an alternating group, or one of the $26$
        sporadics.
  \item \textbf{1981 --- Gorenstein announces completion.} On the
        strength of preprints from many authors, Gorenstein declares
        the classification theorem proved.
  \item \textbf{1983 --- the second-generation project.} Gorenstein,
        Richard Lyons and Ronald Solomon start writing a
        re-organised, streamlined proof. Roughly $13$ planned
        volumes; about $10$ have appeared so far.
  \item \textbf{1989 --- the gap.} Michael Aschbacher and others
        realise that one branch of the case analysis --- the
        \emph{quasithin case} --- was not actually fully proved.
        \emph{Plain version:} ``quasithin'' is jargon for a specific
        narrow strip of the checklist (groups of
        ``characteristic-$2$ type'' with low ``$2$-local rank''),
        and the published proof for that strip had holes. Most
        people thought it would close quickly. It did not.
  \item \textbf{1992 --- Gorenstein dies}, with the proof
        not yet definitively complete, though widely accepted.
  \item \textbf{2004 --- Aschbacher and Smith close the gap.}
        Michael Aschbacher and Stephen Smith publish a two-volume,
        $1{,}221$-page monograph that finally resolves the
        quasithin case. The classification of finite simple groups
        is now universally accepted as proved.
  \item \textbf{2008 onwards --- the second-generation proof rolls
        on.} The Gorenstein--Lyons--Solomon series continues to
        appear, giving a more unified and shorter (still $\sim
        5{,}000$ pages) account.
\end{itemize}

\section*{How did they know $J_4$ was the last one?}

In plain terms: by $1980$, every branch of the giant case analysis
had been resolved.

The strategy from $1955$ onwards was a checklist. Each row of the
checklist was a possible ``shape'' for the centraliser of an
order-$2$ element in a finite simple group. Brauer--Fowler had
guaranteed that each row corresponds to at most finitely many simple
groups, so the checklist had a finite number of rows in principle,
even if nobody knew the rows in advance.

For each row, only four outcomes are possible:

\begin{enumerate}[leftmargin=*, itemsep=2pt]
  \item the row turns out to match a cyclic group $\Z/p\Z$
        (the abelian simple groups);
  \item the row matches an alternating group $A_n$;
  \item the row matches a group of Lie type
        (e.g.\ $\operatorname{PSL}_n(\mathbb{F}_q)$);
  \item the row matches none of the above, in which case a new
        sporadic group has to be constructed to fill it.
\end{enumerate}

The four cases are exhaustive --- this is the content of the
classification theorem. Outcome $(4)$ is the \emph{only} way a new
sporadic enters the list. Each of the $26$ sporadics entered
exactly this way.

By August $1980$:
\begin{itemize}[leftmargin=*, itemsep=2pt]
  \item Every row of the checklist had been worked through.
  \item Every row had been matched to one of the four cases.
  \item No row produced outcome $(4)$ that was not already accounted
        for by one of the $26$.
\end{itemize}

That is what ``$J_4$ is the last one'' actually means: there are no
more rows of the checklist that produce a previously unknown group.
The catch is that ``every row of the checklist has been worked
through'' is itself a theorem, and theorems can have gaps --- the
quasithin gap above is the famous example. But since $2004$, that
gap is closed, and there are no other suspected holes.

So the question ``could there be a 27th?'' reduces to ``could the
proof be wrong?'' That is not zero, but it is a different kind of
question from ``has anyone gone looking?''.

\section*{John Conway and the problem}

John Horton Conway (1937--2020) was a Cambridge lecturer with a
reputation, by his late twenties, for cleverness without a flagship
result. He was conscious of it. The Leech lattice gave him his
problem, and once he had it he kept finding new ways back into the
sporadic-group world for the next half-century.

\subsection*{The Leech lattice (1965)}

In $1965$, John Leech constructed an extraordinarily symmetric lattice
$\Lambda_{24}$ in $24$-dimensional space, arising from the binary
Golay code. It gives the densest known sphere packing in
$24$ dimensions (proved optimal only in $2017$, by Viazovska and
collaborators). Leech suspected its symmetry group
$\operatorname{Aut}(\Lambda_{24})$ was enormous and probably a new
sporadic group, but he could not pin down its order, and he could
not find a group theorist willing to take the problem on.

\subsection*{Conway locks himself in (1968)}

Conway took the bait. The story --- which he told often, with small
variations --- is that he told his wife he would spend twelve hours
every weekday and the full weekend on the problem, ``for as long as
it takes,'' starting on a Saturday. He emerged on the same Saturday
night. The full structure had fallen out in roughly twelve hours.

He had computed
\[
|\operatorname{Aut}(\Lambda_{24})| \;=\; |\mathit{Co}_0| \;=\;
2^{22} \cdot 3^9 \cdot 5^4 \cdot 7^2 \cdot 11 \cdot 13 \cdot 23 ,
\]
about $8.3 \times 10^{18}$, and identified three new sporadic simple
groups in one stroke:
\begin{itemize}[leftmargin=*, itemsep=2pt]
  \item $\mathit{Co}_1 = \mathit{Co}_0 / \{\pm 1\}$;
  \item $\mathit{Co}_2$, the stabiliser of a type-$2$ vector in
        $\Lambda_{24}$;
  \item $\mathit{Co}_3$, the stabiliser of a type-$3$ vector.
\end{itemize}
This is the only time in the history of the classification that
multiple sporadics were uncovered in a single calculation by a
single person. The paper, ``A perfect group of order
$8{,}315{,}553{,}613{,}086{,}720{,}000$ and the sporadic simple
groups'' (1969), made his name.

\subsection*{The Monster, McKay, and moonshine (1973--1979)}

\begin{itemize}[leftmargin=*, itemsep=2pt]
  \item \textbf{1973.} Fischer and Griess independently predict the
        Monster from involution-centraliser arguments. Conway
        becomes one of the people computing what such a group would
        have to look like before anyone has constructed it.
  \item \textbf{1978.} John McKay, reading a number-theory paper on
        the train, notices $196{,}884 = 196{,}883 + 1$ --- the first
        coefficient of the $j$-function written as a sum of Monster
        irreducible-representation dimensions. He mails Conway.
        Conway's first reaction, by his own account, is that it must
        be a coincidence. He checks the next coefficient, then the
        next, and very quickly stops believing in coincidence.
  \item \textbf{1979.} Conway and Simon Norton publish
        \emph{Monstrous Moonshine}, formulating the precise
        conjectures: for each $g \in \M$ there is a specific
        genus-zero modular function $T_g$ (a \emph{Hauptmodul})
        whose $q$-expansion coefficients are the characters of $g$
        on a hypothetical graded $\M$-module. They tabulate $194$
        such functions by hand.
  \item Conway coins the term ``moonshine'' for the phenomenon, in
        the sense of \emph{something so wild it must be nonsense}.
        He kept the name on purpose, even after it stopped looking
        like nonsense.
\end{itemize}

\subsection*{Construction and aftermath (1980--1992)}

\begin{itemize}[leftmargin=*, itemsep=2pt]
  \item \textbf{1980.} Griess constructs $\M$ explicitly as the
        automorphism group of the $196{,}884$-dimensional Griess
        algebra. In the same year, $J_4$ is constructed by a team
        including Conway, Parker, Norton and Benson; the
        classification's last sporadic falls into place.
  \item \textbf{1985.} The \emph{ATLAS of Finite Groups} appears:
        Conway, Curtis, Norton, Parker and Wilson, a single
        oversized volume tabulating character tables, maximal
        subgroups and presentations for every finite simple group up
        to a tractable size, including all $26$ sporadics. It is
        still the reference today.
  \item \textbf{1992.} Borcherds, building on work of Frenkel,
        Lepowsky and Meurman, proves the Conway--Norton conjectures
        by constructing the moonshine module $V^{\natural}$ as a
        vertex operator algebra. He is awarded the Fields Medal in
        $1998$ for this work.
\end{itemize}

\subsection*{The rest of his life}

Conway worked on sporadic groups, geometry of numbers, knot theory,
and recreational mathematics until his death from COVID-19 on
\textbf{11 April 2020}. He is more widely known for the Game of Life,
the surreal numbers, the free will theorem, and the look-and-say
sequence --- but he repeatedly said that the Conway groups and
moonshine were the work he was proudest of. The Monster, he liked
to say, was the only group large enough to make him feel small.

\subsection*{Compact timeline}

\begin{itemize}[leftmargin=*, itemsep=0pt]
  \item \textbf{1937.} Conway born in Liverpool.
  \item \textbf{1964.} PhD, Cambridge (number theory, under Davenport).
  \item \textbf{1965.} Leech constructs $\Lambda_{24}$.
  \item \textbf{1968.} Conway computes
        $\operatorname{Aut}(\Lambda_{24})$ in a single weekend;
        $\mathit{Co}_1, \mathit{Co}_2, \mathit{Co}_3$ are born.
  \item \textbf{1973.} Fischer and Griess predict $\M$.
  \item \textbf{1978.} McKay's observation; Conway is convinced.
  \item \textbf{1979.} Conway--Norton, \emph{Monstrous Moonshine}.
  \item \textbf{1980.} Griess constructs $\M$; $J_4$ constructed;
        list of sporadics closes at $26$.
  \item \textbf{1985.} \emph{ATLAS of Finite Groups} published.
  \item \textbf{1986.} Conway moves from Cambridge to Princeton.
  \item \textbf{1992.} Borcherds proves Monstrous Moonshine.
  \item \textbf{1998.} Borcherds, Fields Medal.
  \item \textbf{2004.} Aschbacher--Smith close the quasithin gap;
        classification accepted complete.
  \item \textbf{2020.} Conway dies, 11 April.
\end{itemize}

\section*{Why ``monster''?}

\begin{itemize}[leftmargin=*, itemsep=2pt]
  \item It has no small faithful representation --- 196{,}883 dimensions
        is the floor.
  \item It contains 19 of the other 26 sporadic groups as subquotients
        (the \emph{Happy Family}). The remaining six are the
        \emph{pariahs}: $J_1$, $J_3$, $J_4$, $\mathit{Ru}$,
        $\mathit{Ly}$, $\mathit{ON}$.
  \item Its connections to string theory, modular forms, lattice theory
        and the Leech lattice are real but still not fully understood.
\end{itemize}

\section*{References and further reading}

\subsection*{For a non-specialist reader}

\begin{itemize}[leftmargin=*, itemsep=2pt]
  \item Mark Ronan, \emph{Symmetry and the Monster: One of the
        Greatest Quests of Mathematics}, Oxford University Press,
        $2006$. A readable narrative history of the entire
        classification, written for a general audience.
  \item Marcus du Sautoy, \emph{Finding Moonshine: A Mathematician's
        Journey Through Symmetry}, Fourth Estate, $2008$ (US title:
        \emph{Symmetry}). A working mathematician's memoir-style
        tour of the subject.
  \item Terry Gannon, \emph{Moonshine Beyond the Monster: The Bridge
        Connecting Algebra, Modular Forms and Physics}, Cambridge,
        $2006$. Heavier going but the best single book on the
        moonshine side.
\end{itemize}

\subsection*{Survey and reference works}

\begin{itemize}[leftmargin=*, itemsep=2pt]
  \item Ronald Solomon, ``A Brief History of the Classification of
        the Finite Simple Groups,'' \emph{Bulletin of the AMS}
        \textbf{38} (2001), 315--352. The standard concise survey
        by one of the people who wrote the proof.
  \item J.\ H.\ Conway, R.\ T.\ Curtis, S.\ P.\ Norton, R.\ A.\
        Parker, R.\ A.\ Wilson, \emph{ATLAS of Finite Groups},
        Clarendon Press, $1985$. The reference work: character
        tables, maximal subgroups and presentations of all $26$
        sporadics and the small finite simple groups generally.
  \item D.\ Gorenstein, R.\ Lyons, R.\ Solomon, \emph{The
        Classification of the Finite Simple Groups}, AMS
        Mathematical Surveys and Monographs $40$, multi-volume from
        $1994$ onwards. The second-generation proof.
\end{itemize}

\subsection*{Key original papers}

\begin{itemize}[leftmargin=*, itemsep=2pt]
  \item R.\ Brauer and K.\ A.\ Fowler, ``On groups of even order,''
        \emph{Annals of Mathematics} \textbf{62} (1955), 565--583.
  \item W.\ Feit and J.\ G.\ Thompson, ``Solvability of groups of
        odd order,'' \emph{Pacific Journal of Mathematics}
        \textbf{13} (1963), 775--1029.
  \item J.\ H.\ Conway, ``A perfect group of order
        $8{,}315{,}553{,}613{,}086{,}720{,}000$ and the sporadic
        simple groups,'' \emph{Bulletin of the London Mathematical
        Society} \textbf{1} (1969), 79--88.
  \item J.\ H.\ Conway and S.\ P.\ Norton, ``Monstrous Moonshine,''
        \emph{Bulletin of the London Mathematical Society}
        \textbf{11} (1979), 308--339.
  \item R.\ L.\ Griess Jr., ``The Friendly Giant,''
        \emph{Inventiones Mathematicae} \textbf{69} (1982),
        1--102.
  \item R.\ E.\ Borcherds, ``Monstrous moonshine and monstrous Lie
        superalgebras,'' \emph{Inventiones Mathematicae}
        \textbf{109} (1992), 405--444.
  \item M.\ Aschbacher and S.\ D.\ Smith, \emph{The Classification
        of Quasithin Groups}, AMS Mathematical Surveys and
        Monographs $111$--$112$, $2004$.
\end{itemize}

\subsection*{Lectures and talks}

\begin{itemize}[leftmargin=*, itemsep=2pt]
  \item Richard Borcherds has posted accessible lecture videos on
        moonshine and vertex algebras on YouTube.
  \item Conway gave many filmed lectures on the Monster and the
        Leech lattice; several are on YouTube under
        ``John Conway Monster.''
\end{itemize}

\section*{The mystery}

We have a construction, and we have a proof, but we still lack a
truly conceptual explanation for \emph{why} the Monster exists, or
why it is bound up with modular functions. Borcherds, on proving the
moonshine conjectures: ``I was over the moon\,\ldots''

\end{document}
